\section{Adaptive Execution Engine}

\label{sec:engine}

\subsection{Execution Model}

The execution engine processes flow graphs using a topological executor with support for parallelism, dynamic expansion, and error recovery.

\begin{algorithm}[H]
\caption{Adaptive Flow Execution}
\label{alg:execute}
\begin{algorithmic}[1]
\Require Flow graph $G = (V, E, \Sigma, \tau)$, input $x_0$
\State $\text{ready} \gets \{v \in V : \text{in-degree}(v) = 0\}$
\State $\text{results} \gets \{\}$
\While{$\text{ready} \neq \emptyset$}
    \State $\text{batch} \gets$ independent nodes in $\text{ready}$
    \State Execute $\text{batch}$ in parallel:
    \For{each $v \in \text{batch}$ \textbf{in parallel}}
        \State $x_v \gets \text{GatherInputs}(v, E, \text{results})$
        \State $y_v \gets \text{ExecuteNode}(v, x_v)$
        \If{$y_v$ is error and $v$ has retry policy}
            \State $y_v \gets \text{RetryWithFallback}(v, x_v)$
        \EndIf
        \State $\text{results}[v] \gets y_v$
    \EndFor
    \State Update $\text{ready}$ with newly satisfiable nodes
\EndWhile
\State \Return $\text{results}[\text{end node}]$
\end{algorithmic}
\end{algorithm}

\subsection{Parallel Branch Execution}

Parallel Split nodes create independent execution branches that run concurrently. The Merge node waits for all (or a configurable quorum of) branches to complete:

\begin{definition}[Merge Strategies]
\begin{itemize}[leftmargin=1.1em]
    \item \textbf{All}: Wait for all branches; output is a list of results.
    \item \textbf{First}: Return the first completed branch; cancel others.
    \item \textbf{Quorum($k$)}: Wait for $k$ of $n$ branches; return the $k$ results.
    \item \textbf{Best}: Wait for all; return the result with the highest quality score.
\end{itemize}
\end{definition}

The ``Best'' strategy is particularly useful for LLM calls: generate multiple responses in parallel with different models or temperatures, then select the best one using a quality evaluator.

\subsection{Dynamic Subgraph Expansion}

Certain patterns require dynamic workflow modification at runtime:

\begin{enumerate}[leftmargin=1.4em]
    \item \textbf{Map-Reduce}: When an LLM Call returns a list, the engine can dynamically spawn a subgraph instance per list item, process them in parallel, and reduce the results.
    \item \textbf{Recursive decomposition}: A task decomposition node can generate subtasks, each of which is processed by a copy of a designated subgraph.
    \item \textbf{Tool-triggered expansion}: When a tool call returns structured data requiring further processing, the engine instantiates a handler subgraph.
\end{enumerate}

\begin{algorithm}[H]
\caption{Map-Reduce Expansion}
\label{alg:mapreduce}
\begin{algorithmic}[1]
\Require Input list $[x_1, \ldots, x_n]$, map subgraph $G_m$, reduce node $R$
\State \Comment{Expand: create $n$ instances of map subgraph}
\For{$i = 1, \ldots, n$ \textbf{in parallel}}
    \State $y_i \gets \text{Execute}(G_m, x_i)$
\EndFor
\State \Comment{Reduce: aggregate results}
\State $y \gets R([y_1, \ldots, y_n])$
\State \Return $y$
\end{algorithmic}
\end{algorithm}

\subsection{Retry and Fallback}

Each node can be configured with a retry policy:

\begin{table}[H]
\centering
\small
\begin{tabular}{lcc}
\toprule
\textbf{Strategy} & \textbf{Behavior} & \textbf{Use Case} \\
\midrule
Fixed retry & Same model, $k$ attempts & Transient errors \\
Backoff retry & Exp. backoff, same model & Rate limits \\
Model fallback & Try next model in list & Model outage \\
Temperature retry & Increase temp each try & Format failures \\
Prompt retry & Append error to prompt & Content errors \\
\bottomrule
\end{tabular}
\caption{Retry strategies available per node.}
\label{tab:retry}
\end{table}

\subsection{Cost and Latency Estimation}

Before execution, the engine estimates total cost and latency:

\begin{equation}
\text{Cost}(G) = \sum_{v \in V} \text{Cost}(v) \cdot \mathbb{E}[\text{retries}(v)],
\end{equation}

\begin{equation}
\text{Latency}(G) = \max_{\text{path } P \in G} \sum_{v \in P} \text{Latency}(v) \cdot \mathbb{E}[\text{retries}(v)],
\end{equation}

where the latency is determined by the critical path (longest sequential chain) due to parallel execution of independent branches.

